% ================================================================
%  chapters/parts/ch08-a.tex   （第 8 章 核磁共振方法，分片 a）
%  源：S2（下册）第 8 章，书内 p.525—543 / PDF p.18—35
%  覆盖：8.1 固体高分辨 NMR 技术（MAS / CP-MAS / 高功率 1H 去偶）
%        8.2 分子筛结构的 MAS NMR 研究（29Si / 27Al / 17O / 31P / 47,49Ti）
%  未覆盖（分片 b/c）：8.3 及以后（酸性、吸附分子、原位、积炭）
%  页码说明：下册印刷页 526 未出现在 PDF 中，故书内页码跳号属正常。
% ================================================================

\chapterstart{核磁共振方法}%
{本章导读：本章回答三个问题——固体谱线为什么极宽；MAS、CP、高功率去偶各自消掉哪一项
相互作用；在什么条件下一条固体 NMR 谱才能被读成结构结论。全章沿「固体高分辨技术
（8.1）$\rightarrow$ 分子筛骨架核（8.2）$\rightarrow$ 表面酸性 $\rightarrow$
吸附分子与原位 $\rightarrow$ 积炭」展开，把每一个谱学量都落到一个可检验的结构
判断上。本分片覆盖 8.1 与 8.2，其中 8.1 是后文全部案例共用的方法与条件基础。}
\label{ch:08}

\subsection{固体高分辨核磁共振技术}

\begin{probout}
固体 MAS NMR 测的是\textbf{原子核局域的电子与几何环境}：化学位移给出配位状态与键合
环境，四极耦合参数给出配位对称性，偶极耦合给出核间距与空间连接，弛豫时间给出运动性。
与 XRD 给出长程有序的晶体结构不同，它给出\textbf{所研究核的局部结构和配位信息}，因此对
短程有序、甚至非晶态的催化材料同样有效；二者结合才能得到较完整的结构图像。
\src{S2 §8，书内 p.527 / PDF p.19}
\end{probout}

固体 NMR 的全部困难可以压缩成一条因果链：

\begin{center}
\fontsize{9}{13}\selectfont
固体中分子不能像液体那样快速各向同性翻滚 $\rightarrow$ 各种内部相互作用不被运动平均
$\rightarrow$ \textbf{谱线极宽}（固体线宽常大于 \qty{100}{kHz}，而液体 \ce{H2O} 的线宽
小于 \qty{1}{Hz}）$\rightarrow$ 分辨率与灵敏度同时塌陷 $\rightarrow$ 必须针对\textbf{每一种}
相互作用设计专门的抵消手段 $\rightarrow$ 才得到可归属、可定量的窄线谱
\end{center}

这条链的关键在于：窄化手段不是通用的「去噪」，\textbf{每一种手段只抵消链上的一到两项}。
把来源与对策一一对应起来，才能判断一张谱在什么条件下可信（表~\ref{tab:narrow-map}）。

\begin{table}[htbp]
\centering
\caption{固体谱线增宽来源与对应窄化手段的一一对应。读者应看第三列：任何一张谱都要说清
它用了哪几种手段，未消除的那一项就是该谱的固有局限。\src{据 S2 §8.1，书内 p.527—531 / PDF p.19—23 归纳。}}
\label{tab:narrow-map}
{\srcsize
\begin{tabularx}{\columnwidth}{@{}lXX@{}}
\toprule
增宽来源 & 表现与量级 & 对应窄化手段 \\
\midrule
同核/异核偶极--偶极 & 与 $\gamma_1\gamma_2/r^3$ 成正比；$^{1}$H、$^{19}$F 同核偶极可达 \qty{100}{kHz} & MAS（空间部分）+ 高功率异核去偶、同核多脉冲（自旋部分） \\
化学位移各向异性（CSA） & 粉末线型宽达数十至数百 ppm & MAS；静态谱反而用来直接读张量主值 \\
四极相互作用（$I>1/2$） & 一阶极大；中心跃迁只余二阶，二阶增宽 $\propto 1/B_0$ & 一阶由 MAS 消除；二阶需高场、MQ MAS 或 DOR/DAS \\
标量（$J$）耦合 & 固体中通常较小 & 去偶、脉冲序列选择 \\
弛豫过快（$T_2$ 短） & 线宽 $\approx 1/(\pi T_2)$ & 降低温度、改善样品均匀性；无法靠脉冲消除 \\
\bottomrule
\end{tabularx}}
\end{table}
\emph{核对说明}：已对照原页面图像（书内 p.531 / PDF p.23）。该页（§8.1.4 高功率
$^{1}$H 去偶）印：同核偶极作用很大时，如 $^{1}$H 和 $^{19}$F 核的同核偶极作用（HD）为
100 kHz，只有满足快速魔角旋转条件，即 $\omega R \gg |HD|_{\max}$，
才可能消除同核偶极作用；一般探头的魔角旋转速度难以达到 111 kHz 超高速 MAS 探头的转速，所以 MAS
技术通常无法消除同核偶极作用；异核去偶（CW 与 TPPM）消除的是 $^{1}$H 对 $^{13}$C
等异核的偶合。本表「同核偶极增宽」一行据该页。化学位移各向异性、四极相互作用、标量 $J$ 耦合、
弛豫过快四行的来源页为 §8.1.3（书内 p.527--531 / PDF p.19--23），
\textbf{尚未逐页核对}。

\begin{principle}
\textbf{增宽的五类来源。}（1）\textbf{核的偶极--偶极相互作用}：同核与异核皆有，大小取决于
核磁矩与核间距；固体中核运动受限，偶极作用比流体样品强，$^{1}$H 等磁矩大的丰核尤其强，
是固体谱增宽的主要因素之一。（2）\textbf{化学位移各向异性}：分子对外磁场取向不同时核外
磁屏蔽与共振频率不同；溶液中各向同性快速运动把它平均成单一值，固体中则表现为宽线，
球对称、轴对称与低对称分子的宽线峰形各不相同。（3）\textbf{四极相互作用}：$I>1/2$ 的核
都存在，溶液中快速翻转把它平均掉，固体中使谱线大大加宽。（4）自旋--自旋标量耦合；
（5）自旋--自旋弛豫时间过短。\textbf{自旋 $I=1/2$ 的核只有（1）（2）两类}；四极核除
（1）（2）外一般以（3）为主。\src{S2 §8.1，书内 p.527—528 / PDF p.19—20}
\end{principle}

\begin{measure}
\textbf{仪器与实验条件。}固体谱仪必须配魔角旋转（MAS）探头，并需要高射频功率
（约 \qty{1000}{W}）以获得短脉冲和宽激发带宽，而液体谱通常只需约 \qty{100}{W}。
转速决定能平均掉多强的相互作用：常见探头为几千赫至十几千赫，可消除 CSA，但只能部分
消除偶极作用；2.5 mm 转子探头可达 \qty{30}{kHz} 以上，最新商业化的 0.7 mm 转子可达
\qty{111}{kHz}。\src{S2 §8.1、§8.1.2，书内 p.528、p.530 / PDF p.20、p.22}
\end{measure}

按自旋量子数，可观测核分为三类，这决定了后面每个核种能问什么问题：$I=1/2$ 核
（$^{1}$H、$^{13}$C、$^{29}$Si、$^{15}$N、$^{31}$P、$^{129}$Xe、$^{195}$Pt 等）；
半整数四极核 $I=3/2,\,5/2,\,7/2,\,9/2$（$^{11}$B、$^{23}$Na、$^{7}$Li、$^{27}$Al、
$^{51}$V、$^{93}$Nb 等）；整数四极核 $I=1,\,3$（$^{2}$H、$^{14}$N、$^{6}$Li、$^{10}$B 等）。
\src{S2 §8.1，书内 p.528 / PDF p.20}

\subsubsection{固体 NMR 的发展过程}

整条发展线索只有一个目标：提高分辨率与灵敏度、建立核与核之间的关联。1958 年 Andrew 与
Lowe 提出魔角旋转（MAS），大大消除了化学位移各向异性与同核偶极作用；1962 年 Hartmann 与
Hahn 提出交叉极化（CP），通过极化转移增强稀核灵敏度；1968 年 Waugh 等实现 WAHUHA 四脉冲
同核去偶；1988 年 Pines 等发展双旋转（DOR）与动态变角旋转（DAS）以消除半整数四极核中心
跃迁的二阶四极作用，因需特殊探头而推广受限；1989 年 REDOR 开创异核间距测量；1994 年
DQ-MAS 建立同核间距测量；1995 年 Frydman 提出多量子魔角旋转（MQ MAS），无需特殊硬件即可
在商用谱仪上消除二阶四极作用，把四极核 NMR 推上一个台阶。本节以下只展开其中三项与分子筛
研究直接相关者。\src{S2 §8.1.1，书内 p.529 / PDF p.21}

\subsubsection{魔角旋转固体核磁共振技术（MAS NMR）}

静态粉末谱直接呈现化学位移各向异性。屏蔽张量的主值为 $\delta_{11}$、$\delta_{22}$、
$\delta_{33}$，由此定义三个可测量量：

\begin{equation}
\label{eq:csatensor}
\begin{aligned}
\delta_{\mathrm{iso}} &= \tfrac{1}{3}\left(\delta_{11}+\delta_{22}+\delta_{33}\right)\\
\Delta\delta &= \delta_{33}-\tfrac{1}{2}\left(\delta_{11}+\delta_{22}\right)\\
\eta &= \left(\delta_{22}-\delta_{11}\right)\big/\left(\delta_{33}-\delta_{\mathrm{iso}}\right)
\end{aligned}
\end{equation}

式中 $\delta_{\mathrm{iso}}$ 为各向同性化学位移，$\Delta\delta$ 为各向异性张量差，
$\eta$ 为不对称因子。\textbf{从静态谱可以直接读出张量的三个主值}；另一个等价途径是拟合
旋转边带，因为边带的包络线就是粉末线型。
（此处已对照原页面图像（书内 p.529 / PDF p.21）：原书逐字印
$\delta_{\mathrm{iso}}=\tfrac{1}{3}(\delta_{11}+\delta_{22}+\delta_{33})$、
$\Delta\delta=\delta_{33}-\tfrac{1}{2}(\delta_{11}+\delta_{22})$、
$\eta=(\delta_{22}-\delta_{11})/(\delta_{33}-\delta_{\mathrm{iso}})$，
$\eta$ 的分母确为 $\delta_{33}-\delta_{\mathrm{iso}}$，三式与上式一致。）
\src{S2 §8.1.2，书内 p.529 / PDF p.21}

使谱线窄化，除高功率 $^{1}$H 去偶外最主要的技术是\textbf{魔角旋转}：把固体样品装在转子中，
绕与外加磁场成 $54.7^\circ$ 的轴高速旋转。旋转样品的磁屏蔽常数写作

\begin{equation}
\label{eq:sigma1}
\begin{aligned}
\sigma_{\mathrm{rot}}&=\sigma_{\mathrm{iso}}
+\frac{1}{2}\left(3\cos^{2}\beta-1\right)\cdot\frac{\delta}{2}
\Bigl[\,3\cos^{2}\theta-1+\eta\sin^{2}\theta\cos 2\phi\,\Bigr]
\end{aligned}
\end{equation}

式中 $\beta$ 为\textbf{样品（转子轴）与外磁场方向 $B_0$ 的夹角}，$\sigma_{\mathrm{iso}}$ 为
各向同性磁屏蔽常数，$\delta$、$\eta$ 反映屏蔽矩阵的各向异性和非对称性，$\theta$、$\phi$
为\textbf{屏蔽环境（分子）坐标系中张量主轴相对转子轴的取向}。
两个角分工不同，不能互换：$\beta$ 由实验者通过转子取向设定，$\theta$、$\phi$ 由分子在
转子中的排列决定、随分子翻滚而变，实验上无法逐个控制。
\textbf{量纲检查}：$\sigma$、$\delta$ 同为磁屏蔽（无量纲量，以 ppm 表示），方括号内两项
皆为「各向异性/各向同性之比」的纯数，故两边量纲一致。
（此处已对照原页面图像（书内 p.529—530 / PDF p.21—22）：原书式 (8-1) 逐字印为
$\sigma_{\mathrm{rot}}=\sigma_{\mathrm{iso}}+\tfrac{1}{2}(3\cos^{2}\beta-1)\cdot
\tfrac{\delta}{2}[\,3(\cos^{2}\theta-1)+\eta\sin^{2}\theta\cos 2\phi\,]$，
即 $\delta$ 确作公因子提到方括号外、$\eta$ 项不带 $1/2$；本处照此改写为公因子形式，
仅把方括号内首项按通行写法作 $3\cos^{2}\theta-1$（原书印 $3(\cos^{2}\theta-1)$，
两者相差常数 2、并不等价；本处按标准式取之，原印法备查）。$\beta=54.7^\circ$ 时整体归零的结论不受影响。）
\src{S2 §8.1.2，书内 p.529—530 / PDF p.21—22}

\begin{principle}
\textbf{魔角条件。}式~\eqref{eq:sigma1} 的后一项是固体特有的化学位移各向异性项，其中只有
$P_2(\cos\beta)=(3\cos^{2}\beta-1)/2$ 这一个因子与转子取向有关。\textbf{令 $P_2$ 为零，
消去的才是转子轴与外场的夹角 $\beta$}（不是局部的 $\theta$）：

\begin{equation}
\label{eq:magic}
\begin{aligned}
&3\cos^{2}\beta-1=0
\;\Longrightarrow\;\cos^{2}\beta=\tfrac{1}{3}\\
&\;\Longrightarrow\;\beta=\arccos\tfrac{1}{\sqrt{3}}=54.7356^\circ
\end{aligned}
\end{equation}

此时 $\sigma_{\mathrm{rot}}=\sigma_{\mathrm{iso}}$，各向异性作用被极大程度消除。\textbf{它
为什么能消掉一阶增宽}：化学位移各向异性、偶极--偶极作用、一阶四极作用的一阶微扰角因子
都只含同一个二阶勒让德多项式 $P_2$；快速旋转下这些含时角因子的\textbf{时间平均}正比于
$P_2(\cos\beta)$。$\beta$ 取 $54.7356^\circ$ 时 $P_2$ 归零，所以\textbf{一项几何条件同时
消掉三类一阶增宽}——这正是「魔角」之名的由来。
$54.7^\circ$ 被称为魔角，样品管在魔角位置上的整体转动即魔角旋转。
\textbf{最常见的读图错误}：$\beta$ 一律\textbf{从 $B_0$（竖直方向）量起}，
不是从水平面量起。把转子轴画成与水平面成 $54.7^\circ$ 时，它与 $B_0$ 只有 $35.3^\circ$，
此时 $P_2(\cos 35.264^\circ)\approx0.5$，各向异性\textbf{只被消去一半、并未归零}，
谱线仍宽。\src{S2 §8.1.2，书内 p.530 / PDF p.22}
\end{principle}

\begin{figure}[htbp]
\centering
% 图源几何自检（可复算）：转子轴方向 = 54.736 deg，自竖直 B_0 量起；
%   端点 (34*sin54.736, 34*cos54.736) = (27.76, 19.63) mm；
%   atan2(27.76, 19.63) = 54.7360 deg  ->  (3cos^2 - 1)/2 = 0
%   若误按 TikZ 极角 (54.736:34) 画（即自水平量起），与 B_0 夹角只有 35.264 deg，
%   P_2 = +0.5，各向异性并未消去。
\begin{tikzpicture}[x=1mm,y=1mm]
  \draw[->,thick] (0,0) -- (0,33)
     node[right,text width=19mm,align=left,font=\fontsize{8}{9.6}\selectfont]{$B_0$：外磁场方向};
  \draw[dashed,gray] (0,0) -- (30,0);
  \draw[gray] (0,11) arc[start angle=90,end angle=54.736,radius=11];
  \node[anchor=south west,font=\fontsize{8}{9.6}\selectfont] at (4.4,11.7)
     {$54.736^\circ$（$P_2=0$）};
  \draw[->,thick] (0,0) -- (27.76,19.63)
     node[anchor=south west,font=\fontsize{8}{9.6}\selectfont] at (28.4,19.0)
     {转子轴（样品管）};
  \begin{scope}[rotate=-35.264]
    \draw[thick,fill=gray!15] (-3.2,-9) rectangle (3.2,9);
  \end{scope}
\end{tikzpicture}
\caption{魔角几何：样品管绕与 $B_0$ 成 $54.736^\circ$ 的轴旋转（$\beta$ 自\textbf{竖直的
$B_0$} 量起，与水平方向成 $35.264^\circ$）。读者应看两处——$\beta$ 由 $P_2(\cos\beta)=0$
定出，而\textbf{三类一阶增宽}（化学位移各向异性、偶极--偶极、一阶四极）\textbf{共用同一个
角因子} $P_2(\cos\beta)=\tfrac{1}{2}(3\cos^{2}\beta-1)$，故只此一个角度就把它们一起归零；
转子轴线与实心样品管长轴重合，图中画成一根轴。本图为几何示意图，
不表示真实探头尺寸与转速方向。\schemamark\figsrc{据 S2 §8.1.2，
书内 p.530 / PDF p.22 改绘。}}
\label{fig:magic-angle}
\end{figure}

\textbf{偶极耦合与核间距。}异核（或同核）两核的偶极耦合常数为

\begin{equation}
\label{eq:dipolar}
b_{12}=-\frac{\mu_0}{4\pi}\cdot\frac{\gamma_1\gamma_2\hbar}{r^{3}},
\qquad
\omega_{\mathrm{D}}(\theta)=b_{12}\cdot\frac{3\cos^{2}\theta-1}{2}
\end{equation}

式中 $\gamma_1$、$\gamma_2$ 为两核的磁旋比，$r$ 为核间距，$\theta$ 为核间矢量与外磁场的夹角。
\textbf{量纲检查}：$\mu_0$ 单位 $\mathrm{H\,m^{-1}}$，$\gamma$ 单位
$\mathrm{rad\,s^{-1}\,T^{-1}}$，$\hbar$ 单位 $\mathrm{J\,s}$，故
$\mu_0\gamma_1\gamma_2\hbar/r^{3}$ 约化后为 $\mathrm{s^{-1}}$（角频率），与 $\omega_{\mathrm{D}}$
一致；除以 $2\pi$ 即得 Hz。两条可直接使用的推论：（i）耦合按 $r^{-3}$ 衰减，核间距增加
10\% 耦合即下降约 25\%，因此偶极耦合是最近的「分子尺」；（ii）角因子 $P_2(\cos\theta)$
在 $\theta=54.736^\circ$ 为零——此处的 $\theta$ 是\textbf{核间矢量与 $B_0$} 的夹角，
与式~\eqref{eq:magic} 的转子轴夹角 $\beta$ 是两个不同的角；MAS 下核间矢量随分子一起
绕转子轴转动，其空间部分的时间平均正比于 $P_2(\cos\beta)$，故转子轴取 $\beta=54.736^\circ$
即可让偶极作用与 CSA 的时间平均\textbf{同时}归零。
\src{S2 §8.1、§8.1.2，书内 p.527、p.530 / PDF p.19、p.22}
\pedset{本处数值算例为教学补充、非原书数据：取 $^{1}$H--$^{1}$H、$r=1.0$~\AA{}，由
式~\eqref{eq:dipolar} 得 $|b|/2\pi\approx\qty{120}{kHz}$，与正文所述 $^{1}$H、$^{19}$F
同核偶极约 \qty{100}{kHz} 同量级，可用作自检。}

\begin{measure}
\textbf{转速够不够，看旋转边带。}MAS 只在「旋转比相互作用快」时把宽线收成一条中心线。
转速不足时，原本连成一片的粉末线型被切成间隔等于转速 $\nu_{\mathrm{r}}$ 的一组
\textbf{旋转边带}（spinning sidebands），中心带只是全部信号的一部分。判断规则有三条：
（i）\textbf{边带间隔等于转速}，可由间隔反读实验转速；（ii）\textbf{边带包络线就是静态
粉末线型}，因此可以用边带强度分布反过来拟合化学位移张量主值；（iii）\textbf{转速足够时
边带强度趋于零}，中心带高度不再随转速增加而上升、线宽也不再收窄。
\src{S2 §8.1.2，书内 p.529—530 / PDF p.21—22}
\end{measure}

\begin{pitfall}
把「中心带面积」当作全部信号来定量，是 MAS 谱最常见的定量错误。只要边带还有可观强度，
按中心带积分就会\textbf{系统性低估}该物种的含量；正确做法是把所有边带一并加和，或把转速
提高到边带可忽略。此外，MAS 只能平均掉偶极哈密顿量的\textbf{空间部分}：当同核偶极作用
很大时（$^{1}$H、$^{19}$F 的同核偶极可达 \qty{100}{kHz}），只有满足 $\omega_{\mathrm{R}}
\gtrsim |H_{\mathrm{D}}|_{\max}$ 的快速旋转条件才能消除它，而一般探头达不到，故
\textbf{MAS 通常并不能消除同核偶极作用}。\src{S2 §8.1.4，书内 p.531 / PDF p.23}
\end{pitfall}

\subsubsection{交叉极化/魔角旋转固体核磁共振技术（CP/MAS NMR）}
\label{sec:cp-mas}

$^{13}$C、$^{15}$N、$^{29}$Si 等稀核丰度低、磁旋比小，检测灵敏度低，且自旋--晶格弛豫
时间 $T_1$ 长，需要很长的弛豫延迟。交叉极化（cross polarization，CP）用丰核（如 $^{1}$H）
的极化去增强稀核信号，条件是两核的射频场满足 Hartmann--Hahn 匹配：

\begin{equation}
\label{eq:hh}
\gamma_{S}B_{1S}=\gamma_{I}B_{1I}
\qquad (I,\,S\ \text{均为}\ I=1/2)
\end{equation}

式中 $S$ 为稀核、$I$ 为丰核。满足匹配即实现极化由丰核向稀核转移，例如 $^{1}$H--$^{13}$C
交叉极化可使 $^{13}$C 信号增强 $\gamma_{\mathrm{H}}/\gamma_{\mathrm{C}}\approx4$ 倍。
CP 脉冲序列中两核同时自旋锁定，其有效的弛豫延迟取决于丰核（$^{1}$H）的
$T_{1\rho}$，而 $^{1}$H 的 $T_{1\rho}$ 往往比 $^{13}$C 的 $T_1$ 短得多，因此信号累加
效率大大提高，对检测季碳尤其有利。当稀核为半整数四极核时，匹配条件变为

\begin{equation}
\label{eq:hhquad}
\nu_{1A}=\left(I+\frac{1}{2}\right)\nu_{1B}
\end{equation}

式中 $A$ 代表 $I=1/2$ 核，$B$ 为四极核。$^{27}$Al 的 $I=5/2$，故静态下 $^{1}$H$\rightarrow
^{27}$Al 的 CP 要求 $\nu_{1\mathrm{H}}=3\nu_{1\mathrm{Al}}$：取 $\nu_{1\mathrm{H}}=
\qty{50}{kHz}$、$\nu_{1\mathrm{Al}}=\qty{16.7}{kHz}$ 即满足匹配。
\src{S2 §8.1.3，书内 p.530 / PDF p.22}

\begin{measure}
\textbf{接触时间是选择性旋钮，也是定量陷阱的来源。}交叉极化的基础是稀核与丰核之间的
偶极--偶极相互作用，不同化学环境的稀核周围丰核的数量与运动状态不同，CP 效率因而不同。
实践中通过改变\textbf{接触时间}（contact time）可以做出选择性：接触时间短时只看到
偶极耦合最强、距 $^{1}$H 最近的核，接触时间长时信号逐渐趋于全谱。因此
\textbf{CP/MAS 谱不能用于定量}，定量必须回到单脉冲（ Bloch 衰减）MAS 谱并配足够长的
弛豫延迟。CP/MAS 与高功率 $^{1}$H 去偶联用，可同时获得高灵敏与高分辨的固体谱。
\src{S2 §8.1.3，书内 p.530 / PDF p.22}
\end{measure}

\subsubsection{高功率 $^{1}$H 去偶技术}

在 MAS 条件下同时对 $^{1}$H 通道施加射频照射，即异核去偶，可进一步窄化 $^{13}$C 等核的
谱线。常用两种方式：\textbf{连续波去偶}（CW decoupling）用最简单的单脉冲连续照射 $^{1}$H，
平均掉 $^{1}$H 对 $^{13}$C 等异核的偶合；\textbf{TPPM}（two-pulse phase-modulated）去偶用
二脉冲相位调制照射 $^{1}$H，是目前应用最广的去偶技术，去偶效率远高于 CW 方法。
\src{S2 §8.1.4，书内 p.531 / PDF p.23}

同核偶极作用需要另一套思路：MAS 平均掉的是偶极哈密顿各向异性的\textbf{空间部分}，
同核去偶脉冲则通过操纵自旋转动消除其\textbf{自旋部分}，只保留各向同性部分。两者互补，
不可相互替代——这也是 8.1.3 小节所述「CP 选择性」与本节「去偶效率」必须分别交代清楚的
原因。

\subsection{分子筛结构的 MAS NMR 研究}

分子筛作为催化剂、催化剂载体与离子交换材料广泛用于工业。沸石分子筛为硅铝酸盐微孔材料，
骨架由 SiO$_{4}$ 与 AlO$_{4}$ 四面体（T 位）构成；SiO$_{4}$ 四面体呈电中性，AlO$_{4}$
四面体带一个负电荷，每个 AlO$_{2}$ 单元需要一个正电荷来平衡，孔道中还含水分子。骨架结构
决定了对不同大小与形状分子的选择催化特征，而针对骨架原子 $^{29}$Si、$^{27}$Al、$^{17}$O、
$^{31}$P 等的 MAS NMR 研究，能非常直接地反映骨架的晶体结构与化学状态。
\src{S2 §8.2，书内 p.531 / PDF p.23}

分子筛骨架组成可写作

\begin{equation}
\label{eq:zeolite}
\mathrm{M}^{x+}_{y/x}\left[\left(\mathrm{AlO}_{2}\right)_{y}
\left(\mathrm{SiO}_{2}\right)_{1-y}\right]\cdot n\,\mathrm{H}_{2}\mathrm{O}
\end{equation}

式中方括号内为骨架，$\mathrm{M}^{x+}$ 为平衡骨架负电荷的阳离子，下标 $y/x$ 为每单位骨架
所需的阳离子数，$n$ 为结晶水数；$x$、$y$ 分别为骨架中 $\mathrm{AlO}_{2}$ 与
$\mathrm{SiO}_{2}$ 单元数。
（此处已对照原页面图像（书内 p.531 / PDF p.23）：原书逐字印
$\mathrm{M}^{x+}_{y/x}[(\mathrm{AlO}_{2})_{y}(\mathrm{SiO}_{2})_{1-y}]\cdot
n\mathrm{H}_{2}\mathrm{O}$，本式与之一致。原文此处关于「每个 $\mathrm{AlO}_{2}$ 单元需要
一个正电荷来平衡」的表述见上一段，与本式下标 $y/x$ 自洽。）
\src{S2 §8.2，书内 p.531 / PDF p.23}

\begin{table}[htbp]
\centering
\caption{分子筛骨架各核种能给出的结构信息与灵敏度难点。读者应看第三、四列：同一份骨架
结构需要多个核种互查，而每个核种都有自己的「看不见」区域。\src{据 S2 §8.2.1--§8.2.5，书内 p.531—539 / PDF p.23—31 归纳。}}
\label{tab:nuclei}
{\srcsize
\begin{tabularx}{\columnwidth}{@{}lXX@{}}
\toprule
核（$I$，天然丰度） & 化学位移范围与典型归属 & 能给出的结构信息与灵敏度难点 \\
\midrule
$^{29}$Si（$1/2$，4.6\%） &
总范围约 500 ppm，分子筛多在 120 ppm 内；Si($n$Al)（$n=0$--4）自高场到低场依次由
Si(0Al) 排到 Si(4Al) &
骨架 Si/Al 比（式~\ref{eq:sial}）、结构不等价 T 位数目与占有率、Si--O--Si 键角。
$T_1$ 长，须足够长弛豫延迟；丰度低、灵敏度低，CP/MAS 可增强但不可定量 \\
$^{27}$Al（$5/2$，100\%） &
总范围约 450 ppm，分子筛多在 100 ppm 内；50--65 ppm 四配位骨架铝，0 ppm 左右六配位
非骨架铝 &
骨架铝/非骨架铝的区分与定量、脱铝程度、铝配位对称性。四极增宽 $\propto1/B_0$，须高场
加高速 MAS；定量须用小扳倒角（$<10^\circ$）强射频脉冲；高度不对称环境可宽到观测不到 \\
$^{17}$O（$5/2$，0.037\%） &
Si--$^{17}$O--Si 与 Si--$^{17}$O--Al 两类氧；3Q MAS 各向同性位移约 50 ppm 与 33 ppm &
骨架 Si/Al 比（由两组分面积）、氧的配位与四极参数（NQCC、$\eta$）。丰度极低，通常必须
$^{17}$O 富集，成本高；同时受 CSA 与四极作用影响 \\
$^{31}$P（$1/2$，100\%） &
AlPO$_{4}$ 与 SAPO 分子筛的 P 位；相变时峰数与强度比改变 &
磷酸铝/硅磷酸铝骨架中 P 的结构位数目与占有率、相变与配位变化。灵敏丰核，$^{31}$P 还常
用作酸性探针（见后文酸性表征部分） \\
$^{47}$Ti、$^{49}$Ti（$5/2$、$7/2$，7.28\%、5.51\%） &
9.4 T 下共振频率 22.547 MHz 与 22.552 MHz；TS-1 约 34 ppm，ETS-10 约 67 ppm，
SrTiO$_{3}$ 为 $-864$ ppm 与 $-1133$ ppm &
钛硅分子筛中 Ti 的配位环境与四极作用大小。丰度不低但四极作用很大、谱线极宽且属低频核，
常需几十万次累加 \\
\bottomrule
\end{tabularx}}
\end{table}
\emph{核对说明}：本表按核种分行，已对照原页面图像的部分为 $^{29}$Si 一行（书内 p.531 /
PDF p.23）：$^{29}$Si 是 $I=1/2$ 核、天然丰度为 4.6\%、化学位移范围在 500
ppm 左右、大多数分子筛的 $^{29}$Si 化学位移均在 120 ppm 范围内。$^{27}$Al、
$^{17}$O、$^{31}$P、$^{47}$Ti/$^{49}$Ti 四行分别落在
§8.2.2--§8.2.5（书内 p.532--539 / PDF p.24--31），
\textbf{尚未逐页核对}。

\subsubsection{$^{29}$Si MAS NMR 研究}

$^{29}$Si 为 $I=1/2$ 核，天然丰度 4.6\%。$^{29}$Si 化学位移取决于分子筛的基本结构，
即 Si($n$Al)（$n=0$--4）与 Si--O--Si 键角；结晶性与磁场强度也影响线宽。
\src{S2 §8.2.1，书内 p.531—532 / PDF p.23—24}

\textbf{低硅铝比分子筛。}含单一结构不等价位的分子筛，其 $^{29}$Si 谱最多可出现五条可分辨
谱峰，对应五种可能的 SiO$_{4}$ 四面体结构，自高场到低场依次为 Si(0Al,4Si)、Si(1Al,3Si)、
Si(2Al,2Si)、Si(3Al,1Si) 与 Si(4Al,0Si)。由于 Loewenstein 规则禁止骨架中出现 Al--O--Al，
可由五种峰的面积 $I$ 直接算出骨架 Si/Al 比：

\begin{equation}
\label{eq:sial}
\frac{\mathrm{Si}}{\mathrm{Al}}=
\frac{\displaystyle\sum_{n=0}^{4}I_{\mathrm{Si}(n\mathrm{Al})}}
{\displaystyle\sum_{n=0}^{4}0.25\,n\,I_{\mathrm{Si}(n\mathrm{Al})}}
\end{equation}

式中分母为铝的摩尔数：每个 Si($n$Al) 贡献 $n$ 个铝，而每个铝由四个 Si 共享，故每条峰
按 $0.25n$ 计权。\textbf{量纲检查}：分子与分母同为峰面积（同一谱的同一标度），比值为纯数，
与 Si/Al 为原子比一致。把式~\eqref{eq:sial} 得到的\textbf{骨架} Si/Al 与传统化学分析
（含骨架铝、非骨架铝和杂质铝）相比，即可估计非骨架铝的量。由于 $^{29}$Si 的 $T_1$ 较长，
必须采用足够长的弛豫延迟才能得到可靠的定量数据；灵敏度不足时可与 $^{29}$Si CP/MAS 联用：
若 CP/MAS 谱中某一谱峰明显增高，说明归属为该峰的是带有 Si--OH 的硅（Q$^{3}$ 类）——
但如前所述，CP/MAS 谱不可用于定量。\src{S2 §8.2.1，书内 p.532—533 / PDF p.24—25}

\textbf{高硅分子筛。}高硅分子筛的 $^{29}$Si 谱只出现 Si(0Al) 及极少量 Si(1Al) 信号，
谱线明显变窄，因为铝含量低使 Si--Al 偶极作用大大减弱；低 Si/Al 样品中大量铝的存在使谱线
宽化，说明 MAS 并未完全消除铝核带来的偶极作用。高硅谱的另一个特点是 Si(0Al) 信号偏向
高场，峰的数目和强度直接反映一个单胞中结构不等价位的数目与含量：八面沸石与 A 型分子筛
只有一个晶格位，只出一条尖锐单峰；丝光沸石有 16T1、16T2、8T3、8T4 四个位，谱线呈
2:1:3 强度比（其中两个位的峰重叠）；$\Omega$ 型分子筛有 24T1 与 12T2 位，出现强度比
2:1 的两个峰，容易归属为 T1 与 T2。注意式~\eqref{eq:sial} \textbf{只适用于具有单一 T 位的
分子筛系统}。\src{S2 §8.2.1，书内 p.533 / PDF p.25}

窄谱线还使 $^{29}$Si MAS NMR 能跟踪骨架的相变与吸附行为：ZSM-5 在室温吸附对二甲苯时，
吸附量为 0.4 分子/单胞时谱图变化不大，仍呈对应单斜晶格 24 个 T 位的 24 条谱线；吸附量增
至 1.6 分子/单胞时谱图完全变为 12 条谱线，表明发生了单斜晶格向正交晶格的转变；在 0.4--1.6
分子/单胞之间，两种晶格以不同比例并存。\src{S2 §8.2.1，书内 p.534 / PDF p.26}

\begin{figure}[htbp]
\centering
\begin{tikzpicture}
\tikzset{declare function={lor(\x,\c,\w)=1/(1+((\x-\c)/\w)^2);}}
\begin{axis}[
  width=\onecolwidth, height=56mm,
  xmin=-42, xmax=42, ymin=-0.4, ymax=7.9,
  xtick={-30,-15,0,15,30}, ytick=\empty,
  xlabel={$\delta$ / ppm}, ylabel={信号强度（任意单位）},
  axis line style={black}, tick align=outside,
  every axis plot/.append style={thick},
  samples=900
]
\addplot[black,domain=-42:42]{0.05+1.00*lor(x,-11,13)+1.00*lor(x,11,13)};
\node[anchor=west,text width=15mm,align=left,font=\fontsize{8}{9.6}\selectfont]
  at (axis cs:23,0.75){静态粉末谱};
\addplot[black,domain=-42:42]{2.00+1.55*lor(x,-11,0.85)+1.55*lor(x,0,0.85)+1.55*lor(x,11,0.85)};
\node[anchor=west,text width=15mm,align=left,font=\fontsize{8}{9.6}\selectfont]
  at (axis cs:23,2.70){MAS};
\addplot[black,domain=-42:42]{4.00+2.60*lor(x,-11,0.85)+2.60*lor(x,0,0.85)+2.60*lor(x,11,0.85)};
\node[anchor=west,text width=15mm,align=left,font=\fontsize{8}{9.6}\selectfont]
  at (axis cs:23,5.20){CP/MAS};
\end{axis}
\end{tikzpicture}
\caption{同一组三条共振线的增宽与窄化对照（自下而上：静态粉末谱、MAS 谱、CP/MAS 谱）。
读者应看三点：静态谱因化学位移各向异性与偶极作用连成一片宽包；MAS 把它收成三条窄线，
但峰高受灵敏度限制；CP/MAS 在同样线宽下信号大幅增强。本图按洛伦兹线型计算的示意图，
不是任何具体样品的实测谱，纵轴为任意单位，三条谱线按竖直方向错开以便比较。
\schemamark\figsrc{据 S2 §8.1.2--§8.1.3，书内 p.529—530 / PDF p.21—22 的原理绘制。}}
\label{fig:mas-compare}
\end{figure}

\subsubsection{$^{27}$Al MAS NMR 研究}

$^{27}$Al 是骨架的另一个关键核，天然丰度 100\%，化学位移总范围约 450 ppm，分子筛多数在
100 ppm 内。$^{27}$Al 是 $I=5/2$ 的四极核：由于固体中的四极相互作用，只有
$(+1/2\leftrightarrow-1/2)$ 中心跃迁可以观测到，该跃迁的线形畸变与位移不受一阶相互作用
影响，仅由二阶相互作用决定；其他允许跃迁因谱线太宽、位移太远而无法直接观测。
\textbf{四极作用引起的增宽与磁场强度成反比}，故高场加快速旋转才能得到高质量的
$^{27}$Al MAS 谱；要得到定量可靠的结果，还须采用扳倒角小于 $10^\circ$ 的强射频脉冲。
同一 Y 型分子筛在 \qty{23.45}{MHz} 与 \qty{104.22}{MHz} 下的谱对比表明，高场获得的谱线
更窄、畸变更小——但即使如此，MAS 也只能大幅减小而\textbf{一般不能完全消除}四极相互作用。
结构完整的沸石分子筛中铝的环境均为 Al(4Si)，四面体铝只在 50--65 ppm 出现单个共振峰；
$^{27}$Al 化学位移对不同配位的铝物种十分敏感，因此可据此区分四配位骨架铝与 0 ppm 附近的
六配位非骨架铝。\src{S2 §8.2.2，书内 p.534—535 / PDF p.26—27}

脱铝过程提供了这套判据的典型用法：结构规整的 Y 型分子筛只在 58 ppm 左右出现一个四配位
骨架铝峰；与 SiCl$_{4}$ 反应后骨架铝峰大大降低，同时在 0 ppm 左右出现八面体非骨架铝峰，
并在 100 ppm 处出现一个很强的峰，归属为高温反应生成的 NaCl 与 AlCl$_{3}$ 进一步生成的
Na[AlCl$_{4}$]；经水洗除去 Na[AlCl$_{4}$] 后 100 ppm 峰消失，但更多非骨架铝沉积在孔道中，
0 ppm 峰增强；再经多次水洗或离子交换可除去孔道中大量非骨架铝。水洗后谱线的宽化说明铝物种
环境的\textbf{不对称性增加}：$^{27}$Al 谱的线宽主要取决于核四极耦合常数
$C_{Q}=e^{2}qQ/h$（$q$ 为电场梯度张量的 $z$ 分量，$eQ$ 为核四极矩）与不对称因子 $\eta$，
高度不对称的环境会使谱线严重增宽，直至一些 $^{27}$Al 信号难以被观测到。
\src{S2 §8.2.2，书内 p.536 / PDF p.28}

\begin{pitfall}
\textbf{「看不见」不等于「不存在」。}这是 $^{27}$Al MAS NMR 最容易出错的一处：当铝落进
高度不对称的配位环境时，二阶四极增宽可把信号摊到极宽的频率范围内，中心带峰高降到噪声
以下而被判为「无铝」。正确做法是先用高场、高速 MAS 与回波（或 MQ MAS）确认信号是否真的
不存在，再下结论；把「未观测到」直接写成「不含该物种」，会把结构结论建立在仪器条件上。
\src{S2 §8.2.2，书内 p.536 / PDF p.28}
\end{pitfall}

\subsubsection{$^{17}$O MAS NMR 研究}

$^{17}$O 与 $^{27}$Al 一样是 $I=5/2$ 的四极核，但天然丰度只有 0.037\%，通常必须做
$^{17}$O 富集。$^{17}$O MAS NMR 谱\textbf{同时受化学位移各向异性与核四极矩的影响}，这一
双重性正是它的信息价值所在。A 型分子筛中 $^{17}$O 只有 Si--$^{17}$O--Al 结构，其谱峰显示
较小的四极作用特征；而高硅 Y 型分子筛中 $^{17}$O 只含 Si--$^{17}$O--Si 结构，却显示一定
的四极相互作用。Si/Al 比为 2.74 的 NaY 分子筛，其 $^{17}$O 静态谱与 MAS 谱都可以由
Si--$^{17}$O--Si 和 Si--$^{17}$O--Al 两个组分模拟；由这两个组分的峰面积即可计算 Si/Al 比，
且这样得到的\textbf{仅仅是骨架中的} Si/Al 比——这是 $^{17}$O 相对于 $^{29}$Si 与化学分析
的独特之处。\src{S2 §8.2.3，书内 p.536 / PDF p.28}

半整数四极核（$^{23}$Na、$^{27}$Al、$^{17}$O 等）的研究兴趣推动了多种消除二阶四极作用的
技术：二维章动谱、变角实验、动态变角（DAS）、双旋转（DOR）与多量子谱等。其中 MQ MAS 比
DOR、DAS 简单得多，不需要特殊硬件，特别适用于自旋--晶格弛豫时间短和四极耦合常数较小的核。
多量子 MAS 实验通常用一个简单的二脉冲序列得到二维谱，通过相位循环选择三量子或五量子
跃迁，得到 3Q MAS 或 5Q MAS 谱。以 $^{17}$O 富集的 Na-ZSM-5 的二维 $^{17}$O 3Q MAS 谱
为例：F2 域投影就是常规一维 $^{17}$O MAS 谱，两个相关峰表明存在两种不同的 $^{17}$O 物种；
F1 域投影出现两个分离峰，可直接读出两种物种的各向同性化学位移分别为 50 ppm
（Si--O--Si）与 33 ppm（Si--O--Al），NQCC 分别为 5.3 MHz 与 3.5 MHz，$\eta$ 为 0.12 与
0.29；再据这些四极参数模拟一维 MAS 谱，得到两种物种的浓度分别为 80\% 与 20\%。MQ MAS 的
主要局限是\textbf{不适宜研究四极相互作用很大的物种}。
\src{S2 §8.2.3，书内 p.537—538 / PDF p.29—30}

\subsubsection{$^{31}$P MAS NMR 研究}

$^{31}$P 的 $I=1/2$，天然丰度 100\%，是很适宜 NMR 研究的灵敏丰核。含磷的磷酸铝分子筛
（如 AlPO$_{4}$-11）及含 P、Si、Al 的 SAPO 类分子筛，除研究其 $^{29}$Si 与 $^{27}$Al 谱外，
还可用 $^{31}$P MAS NMR 研究结构变化。VPI-5 的原位变温 $^{31}$P MAS NMR 谱是一个干净的
例子：室温下 VPI-5 有三个强度比 1:1:1 的 $^{31}$P 峰，对应空间群 P6$_{3}$ 结构中三个
占有率为 1:1:1 的 P 结构位；随温度升高，三个峰最终变为强度比 2:1 的两个峰，对应空间群
P6$_{3}$cm 结构中两个占有率为 2:1 的 P 结构位。峰的数目与强度比直接读出结构位的数目与
占有率，这正是 $^{31}$P MAS NMR 在磷酸铝体系中的主要用法。
\src{S2 §8.2.4，书内 p.538 / PDF p.30}

\subsubsection{$^{47,49}$Ti MAS NMR 研究}

$^{49}$Ti 的 $I=7/2$、天然丰度 5.51\%；$^{47}$Ti 的 $I=5/2$、天然丰度 7.28\%。两者的
丰度都不算低，但有很大的四极相互作用，谱线极宽，且属于低频共振核：在 9.4 T 磁场下
$^{47}$Ti 与 $^{49}$Ti 的共振频率分别为 22.547 MHz 与 22.552 MHz。钛硅分子筛 TS-1 与
ETS-10 是两类受关注的材料，TS-1 为高硅 MFI 型分子筛，骨架中有 0.1\%--2.5\% 的 Si 被 Ti
取代；ETS-10 为大孔分子筛。SrTiO$_{3}$ 的 $^{49}$Ti 与 $^{47}$Ti 分别给出 $-864$ ppm 与
$-1133$ ppm 两条较窄的峰，说明其钛的四极相互作用较小；TS-1 与 ETS-10 则分别在 34 ppm
与 67 ppm 出 $^{49}$Ti 峰，由于存在很强的四极相互作用，谱峰很宽、灵敏度很低，均需几十万
次累加才能得到可用的 Ti MAS NMR 谱。\src{S2 §8.2.5，书内 p.539 / PDF p.31}

\begin{dataflow}
\textbf{从一张 MAS 谱到一条结构结论。}（1）\textbf{确认采集条件}：定量结论要求弛豫延迟
足够长（$^{29}$Si 的 $T_1$ 长，尤其如此），四极核定量还要求小扳倒角脉冲；谱是 CP/MAS 时
只能谈存在性，不能谈比例。（2）\textbf{定标}：用公认参比确定化学位移零点，跨磁场比较时
须用同一参比。（3）\textbf{归属}：先按化学位移分区（如 $^{27}$Al 的 50--65 ppm 四配位 /
0 ppm 六配位），再用已知结构样品校核分界值。（4）\textbf{定量}：分峰拟合求面积，并检查
旋转边带是否已全部计入。（5）\textbf{换算}：把面积转成结构量，如用式~\eqref{eq:sial}
得骨架 Si/Al 比、用 $^{17}$O 两组分面积得骨架 Si/Al 比。（6）\textbf{回查}：把结果与独立
方法对照（见下）。任何一步缺失，都不应把数字写成结构结论。
\src{S2 §8.2.1--§8.2.5，书内 p.531—539 / PDF p.23—31}
\end{dataflow}

\begin{crosscheck}
\textbf{互补与回查。}固体 NMR 给的是局域信息，必须与其它方法互查才能得到完整的结构图像：
（i）与 \textbf{XRD} 互补——XRD 给长程有序的晶体结构与结晶度，NMR 给局部配位；短程有序
或非晶态材料中后者仍有效。（ii）与\textbf{化学分析}互补——式~\eqref{eq:sial} 给出骨架
Si/Al 比，化学分析给出总 Si/Al 比，两者之差即非骨架铝的量，这一差值是判断脱铝程度的
常用依据。（iii）\textbf{多核互查}——$^{29}$Si 看硅的分布与 Si/Al 比，$^{27}$Al 看骨架铝
与非骨架铝的区分，$^{17}$O 给出骨架 Si/Al 比与氧的四极参数，$^{31}$P 用于磷酸铝体系，
$^{47,49}$Ti 用于钛硅材料；同一份样品上若 $^{29}$Si 给出的骨架 Si/Al 比与 $^{17}$O 不
一致，应先怀疑定量条件（弛豫延迟、边带、脉冲扳倒角），而不是直接采信任一方。
（iv）与\textbf{红外、NH$_3$ 程序升温脱附}等在酸性表征上互查（见后文酸性表征部分）。
\src{S2 §8.1--§8.2，书内 p.527—539 / PDF p.19—31}
\end{crosscheck}
